\documentclass[11pt]{article}
\usepackage[a4paper,margin=25mm]{geometry}
\usepackage{amsmath,amssymb,amsthm,hyperref}
\newtheorem{lemma}{Lemma}\newtheorem{theorem}{Theorem}
\title{A localized mass bound and weighted maximum principle}
\author{Conditional analytic theorem; three independent mathematical reviews completed}
\date{30 September 2026}
\begin{document}\maketitle
This note isolates the measure-theoretic final step of the proposed
endpoint-profile improvement. It does not assert that all its kernel
comparison hypotheses have already been proved for the dice tensor.
All intervals below are fixed compact subintervals of $(0,\infty)$.
Write $b=\sqrt{\tau}$, where $\tau=C_0\log m/m$, and let $J_0$ lie
strictly inside the interior of $J=[a_-,a_+]$.

We use positive kernels $K_t(x,a)$, for $t=\tau,c\tau$ with fixed $c>1$,
which on $J$ satisfy fixed Gaussian upper bounds
\[
 K_t(x,a)\le Cb^{-1}e^{-|x-a|^2/(Cb^2)},
 \qquad K_\tau(x,a)\ge c_1b^{-1}\quad (|x-a|\le b).
\]
The argument is unchanged if all distances are in $\sqrt{x}$, since
that coordinate is uniformly bi-Lipschitz on $J$.
All measures have mass bounded by a fixed constant. Kernel integrals
outside $J$, for centers at distance $\delta$ from its boundary, are
assumed to be at most
$C m^C b^{-1}\exp[-c\delta^2/b^2]$; the same estimate is assumed after
multiplication by any auxiliary functions used below. This can instead
be supplied with exponential-moment weights.
Suppose a reference measure $\beta$ has a density bounded above and
below by positive constants near $J$, so
$c\le\beta(K_\tau(\cdot,a))\le C$ for centers separated from the
boundary by much more than $b$.

\begin{lemma}[Absorbing a broader localized kernel]\label{mass}
Suppose a positive measure $\mu$ satisfies, uniformly for $a\in J$,
\[
 |\mu(K_\tau(\cdot,a))-\beta(K_\tau(\cdot,a))|
 \le \epsilon_m\mu(K_{c\tau}(\cdot,a))+m^{-B},
 \qquad \epsilon_m\longrightarrow0,
\]
where $B$ can be any sufficiently large fixed constant.
Then uniformly for $a\in J_0$,
\[
 \mu([a-b,a+b])\le Cb,\qquad
 c\le\mu(K_\tau(\cdot,a))\le C.
\]
Consequently, for $a\in J_0$ and $h\ge b$ whose $h$-neighborhood
remains in a fixed interior subinterval of $J$,
\[
 \int_{|x-a|\ge h}K_\tau(x,a)\,d\mu(x)
 \le C e^{-c h^2/b^2}+O(m^{-B}).
\]
\end{lemma}
\begin{proof}
Let $d(a)=\min(a-a_-,a_+-a)$, and use the weight $\phi(a)=d(a)^2$.
Set
$Q=\sup_{a\in J}\phi(a)\mu([a-b,a+b])/b$.
It suffices to treat $Q\ge1$. Choose an approximate maximizer $a$
with value at least $Q/2$. The total-mass bound gives
$\phi(a)\ge cQb\ge cb$, so its boundary distance
$\delta=d(a)$ satisfies $\delta\ge c\sqrt b\gg b$.
On $|x-a|\le\delta/2$, the weights $\phi(x)$ and $\phi(a)$ are
comparable. Partition that neighborhood into intervals of length $b$;
the Gaussian upper bound and the definition of $Q$ give
\[
 \mu(K_{c\tau}(\cdot,a))\le C Q/\phi(a)
       +C m^Cb^{-1}e^{-c\delta^2/b^2}.
\]
The second term is superpolynomially small because $\delta^2\ge cb$.
The lower kernel bound gives
$\mu(K_\tau(\cdot,a))\ge cQ/\phi(a)$.
Multiply the comparison hypothesis by $\phi(a)$ to obtain
$Q\le C+C\epsilon_m Q+o(1)$. Absorption proves $Q\le C$.
This gives the local interval bound on every fixed interior subinterval.
Summing Gaussian tails over the same length-$b$ partition proves the
last assertion and the uniform upper bound on broader kernel masses.
The comparison hypothesis and the positive reference integral now give
the positive lower bound.
\end{proof}

\begin{theorem}[Weighted maximum principle]\label{maximum}
Suppose there are finitely many marked row systems with positive
measures $\mu_i$, scalar coordinates $x_i$, and functions
$h_i=z_i-x_i$, where $z_i$ is nondecreasing as a function of the row
coordinate $x_i$. Flat rows are allowed, with the corresponding order
interpretation. Suppose the local-mass and kernel-mass conclusions of
Lemma~\ref{mass} hold uniformly for these measures for all centers in
$J$. This follows by applying that lemma on a slightly larger fixed
interval, with $J$ as its interior interval.
Assume a prior polynomial bound
\[
 |h_i|\le R_m\le m^{-\alpha},\qquad \alpha>0,
\]
on $J$, after weakening a power slightly to absorb logarithmic factors.
Let $H_i(x)$ control the row range occurring in the marked error.
For the weight $\phi(x)=d(x)$, assume that,
if $D=\max_{i,\text{rows}}\phi(x_i)|h_i|$ and $a\in J$, then
\[
 H_i(x)\le C D/\phi(a)
 \quad\text{whenever }|x_i-a|\le d(a)/2.
 \tag{W}
\]
It is enough that (W) hold with a fixed change in the fraction $1/2$.
Finally assume the localized marked comparison
\[
 \left|\int h_i(x)K_\tau(x,a)\,d\mu_i(x)\right|
 \le {C\over m}\int H_i(x)^2
       \left(1+b^{-1}+H_i(x)b^{-2}\right)
       K_{c\tau}(x,a)\,d\mu_i(x)+m^{-B},
 \tag{M}
\]
with the stated negligible tail bounds. Then
\[
 \sup_{i,\ x_i\in J_0}|h_i(x_i)|
       \le C_{J_0,J}\sqrt{\log m/m}.
\]
\end{theorem}
\begin{proof}
Suppose $D>C_*b$, and take a maximizing marked row, with coordinate
$x_0$, signed deviation $h_i(x_0)$, absolute deviation $\Delta$, and
boundary distance $\delta=\phi(x_0)$.
Then $\Delta\le R_m$ and
\[
 \delta=D/\Delta\ge C_*b/R_m,
 \qquad \delta/b\ge C_*/R_m\longrightarrow\infty.
\]
The ratio grows polynomially, so Gaussian tails beyond $\delta/3$
are superpolynomially small. We only shift the test center by a fixed
multiple of $b$, not by the potentially much larger $\Delta$.

Choose fixed constants $B_0$ and $A_0=2B_0$, in that order, with $B_0$
sufficiently large for the Gaussian tails below. If
$h_i(x_0)=\Delta>0$, center the kernel at $a=x_0+A_0b$.
For $|x-a|\le B_0b$, monotonicity gives
\[
 h_i(x)\ge\Delta-(x-x_0)
       \ge\Delta-(A_0+B_0)b\ge\Delta/2,
\]
provided $C_*$ is large, since $\Delta=D/\delta\ge c_JC_*b$.
For a negative deviation center at $a=x_0-A_0b$ and reverse signs.
All these intervals remain inside $J$, and their cutoff weights differ
by a factor $1+o(1)$ from $\delta$.
Throughout the larger $\delta/3$ neighborhood, weighted maximality
gives $|h_i|\le C\Delta$, and (W) gives $H_i\le C\Delta$.
By Lemma~\ref{mass}, the total kernel mass is at least a constant,
whereas the mass outside $|x-a|\le B_0b$ is at most
$C\exp(-cB_0^2)+o(1)$. Taking $B_0$ large therefore yields
\[
 \left|\int h_i K_\tau(\cdot,a)\,d\mu_i\right|
 \ge c\Delta.
\]
The far part outside the $\delta/3$ neighborhood is negligible by
the preceding polynomial lower bound on $\delta/b$.
The right side of (M), using the bounded broader kernel mass, is at most
\[
 {C\Delta^2\over m}\left(1+b^{-1}+\Delta b^{-2}\right)+m^{-B}
 =o(\Delta),
\]
because $\Delta\le R_m\to0$, $mb^2=C_0\log m$, and $mb\to\infty$.
This contradiction proves $D\le C_*b$. Since $\phi$ is bounded below
on $J_0$, the conclusion follows.
\end{proof}

\paragraph{Tensor bookkeeping still required.}
For the actual marked coordinates, $x_i$ depends on the deleted column.
The full-mean deviations satisfy the exact relation
$\ell u_i-x_i=(\ell/M)(m u_i-s)$.
The coordinate differences between deleted columns are $O(R_m/m)$
on a compact endpoint range. Thus (W) follows from the weighted maximum
of the full-mean deviations, since the maximizing row is much farther
than $b$ (and than the $O(R_m/m)$ coordinate discrepancy) from the
boundary. This last conversion should be written
explicitly in the application, together with spectral truncation and
global-tail estimates supplying the two comparison hypotheses.
\end{document}
